% Options for packages loaded elsewhere
\PassOptionsToPackage{unicode}{hyperref}
\PassOptionsToPackage{hyphens}{url}
\PassOptionsToPackage{dvipsnames,svgnames,x11names}{xcolor}
%
\documentclass[
  11pt,
]{article}
\usepackage{amsmath,amssymb}
\usepackage{iftex}
\ifPDFTeX
  \usepackage[T1]{fontenc}
  \usepackage[utf8]{inputenc}
  \usepackage{textcomp} % provide euro and other symbols
\else % if luatex or xetex
  \usepackage{unicode-math} % this also loads fontspec
  \defaultfontfeatures{Scale=MatchLowercase}
  \defaultfontfeatures[\rmfamily]{Ligatures=TeX,Scale=1}
\fi
\usepackage{lmodern}
\ifPDFTeX\else
  % xetex/luatex font selection
\fi
% Use upquote if available, for straight quotes in verbatim environments
\IfFileExists{upquote.sty}{\usepackage{upquote}}{}
\IfFileExists{microtype.sty}{% use microtype if available
  \usepackage[]{microtype}
  \UseMicrotypeSet[protrusion]{basicmath} % disable protrusion for tt fonts
}{}
\makeatletter
\@ifundefined{KOMAClassName}{% if non-KOMA class
  \IfFileExists{parskip.sty}{%
    \usepackage{parskip}
  }{% else
    \setlength{\parindent}{0pt}
    \setlength{\parskip}{6pt plus 2pt minus 1pt}}
}{% if KOMA class
  \KOMAoptions{parskip=half}}
\makeatother
\usepackage{xcolor}
\usepackage[margin=1in]{geometry}
\usepackage{longtable,booktabs,array}
\usepackage{calc} % for calculating minipage widths
% Correct order of tables after \paragraph or \subparagraph
\usepackage{etoolbox}
\makeatletter
\patchcmd\longtable{\par}{\if@noskipsec\mbox{}\fi\par}{}{}
\makeatother
% Allow footnotes in longtable head/foot
\IfFileExists{footnotehyper.sty}{\usepackage{footnotehyper}}{\usepackage{footnote}}
\makesavenoteenv{longtable}
\setlength{\emergencystretch}{3em} % prevent overfull lines
\providecommand{\tightlist}{%
  \setlength{\itemsep}{0pt}\setlength{\parskip}{0pt}}
\setcounter{secnumdepth}{-\maxdimen} % remove section numbering
\ifLuaTeX
  \usepackage{selnolig}  % disable illegal ligatures
\fi
\usepackage{bookmark}
\IfFileExists{xurl.sty}{\usepackage{xurl}}{} % add URL line breaks if available
\urlstyle{same}
\hypersetup{
  pdftitle={Slot grammar in the normalized terminal-triple branch},
  pdfauthor={Aafi Khan Malek (Independent Researcher, Ahmedabad, India)},
  colorlinks=true,
  linkcolor={Maroon},
  filecolor={Maroon},
  citecolor={Blue},
  urlcolor={Blue},
  pdfcreator={LaTeX via pandoc}}

\title{Slot grammar in the normalized terminal-triple branch}
\usepackage{etoolbox}
\makeatletter
\providecommand{\subtitle}[1]{% add subtitle to \maketitle
  \apptocmd{\@title}{\par {\large #1 \par}}{}{}
}
\makeatother
\subtitle{Twenty exact branch exclusions for nine separated subset sums}
\author{Aafi Khan Malek (Independent Researcher, Ahmedabad, India)}
\date{November 2026}

\begin{document}
\maketitle

\section{Slot grammar in the normalized terminal-triple
branch}\label{slot-grammar-in-the-normalized-terminal-triple-branch}

Assume a hypothetical nine-weight configuration in the established
floor-148 slice, in its close-triple representative \[
Y=x_7=X-1,\qquad X=x_8\in[148,149),\qquad M=x_9=X+1 .
\] The unique short-block normalization preserves every labeled subset
comparison, so all earlier central rank restrictions remain valid for
this representative. In particular \([158,355]\) has every consecutive
real gap strictly below 2, and the eight prescribed length-148
coordinate-9 edges have the known top-bit patterns. We use the
established convention \(D_j=s_{256+j}-s_{257-j}\) for the \(j\)-th
ordered positive absolute full-sign magnitude (one per complementary
pair). Writing \(\delta=M-149=X-148\), the known central envelope is \[
2j-1\le D_j\le2j+\delta,\quad 1\le j\le64.
\] Adjacent \(D_j\)'s are at least 2 apart: both complementary ranks
move apart by at least one real unit. For the indices used below, they
are also strictly less than 4 apart, since
\(D_{j+1}-D_j=2(s_{257+j}-s_{256+j})\) by complement and these
consecutive ranks are within the central range.

Positive first-six subset differences are at least 3 apart from zero,
and avoid \((X-3,X+3)\), by the normalized top differences 1,2 and the
middle band \(X-2,\ldots,X+2\). In particular, first-six full absolute
signed magnitudes are distinct one per complementary pair, and are
spaced at least 6. The spacing follows by applying the same
complementary-rank identity to the ordered 64 subset sums of these six
coordinates, with unit spacing replaced by 3. Any nonzero signed
coefficient word in \(\{-1,0,1\}^6\) is a difference of two first-six
subsets, so it too must avoid that open forbidden interval in positive
orientation.

\subsection{1. Complete trio partition}\label{complete-trio-partition}

In a representative of each full nine-sign complement pair choose at
least two top signs plus. This selects exactly one raw sign word per
complementary pair. For each of the 64 first-six sign words, let \(b\)
be its signed first-six sum. The raw top-three contributions for two
pluses are \(X-2,X,X+2\), giving full raw sums \(b+X-2,b+X,b+X+2\). Here
\(|b+X|\ge3\): if \(b<0\) this follows from the middle forbidden
interval, and if \(b>0\) from \(X>0,b\ge3\). Thus the three raw signs
are the same and taking absolute values produces a trio with gap exactly
2 at both steps. No fourth magnitude can intervene, by minimum 2
spacing.

In positive full orientation the ascending top patterns within a
two-plus trio are \(++-\), \(+-+\), \(-++\). Within a one-plus trio they
are \(+--\), \(-+-\), \(--+\); these arise by complementing a negative
raw trio. The all-plus raw form gives one other magnitude \(|b+3X|\),
whose positive top orientation is \(+++\) or \(---\). Thus the entire
list is a disjoint tiling by 64 complete trios and 64 other singleton
magnitudes. A singleton here means its block type, not a claim about
nearby numerical gaps. All identities concern actual labeled signed
words: if two distinct complement pairs had equal absolute sum, their
positive orientations would give identical complementary rank gaps,
contrary to strict rank separation. No ambiguity is introduced by taking
positive orientations.

\subsection{2. All four prescribed pairs are two-plus
mirrors}\label{all-four-prescribed-pairs-are-two-plus-mirrors}

For \(q=50,53,60,63\), the known formulas are \[
D_q=v_q+1,\qquad D_{q+1}=2X+1-w_{q+1},
\] where \(v_q,w_{q+1}\) are positive oriented first-seven full forms.
The top signs for \(D_q\) are \((v_q(7),-,+)\), and for \(D_{q+1}\) are
\((-w_{q+1}(7),+,+)\). If \(v_q(7)=+\), \(D_q\) is middle in a two-plus
trio and \(D_{q+1}\) must be its last member. Hence \(w_{q+1}(7)=+\) and
the first-six word of \(w\) is the negative of the first-six word of
\(v\). Conversely if \(w_{q+1}(7)=+\), its trio predecessor at \(q\) has
top pattern \(+-+\), giving \(v_q(7)=+\). Thus the seventh signs match.
If both are minus, \(D_q\) ends a one-plus trio and \(D_{q+1}\) is an
all-plus positive singleton.

For \(q=53,63\), the both-minus option is impossible: the first term of
its one-plus trio at \(q-2=51,61\) would have top signs \(+--\), but
these preceding forced singleton-or-last-trio positions have top
coordinates 8 and 9 both plus by their displayed formula. Therefore
\(v,w\) are both plus at \(53,63\).

Call a singleton of index in the range used here \emph{low} when
\(D=3X-a\) for a positive first-six magnitude \(a<3X\), and \emph{high}
when \(D=a-3X\) with \(a>3X\). These exhaust the possibilities here: a
raw magnitude \(3X+a\) is far above \(D_{64}<129\). For a two-plus trio
center \(D_c\) in this range the corresponding positive first-six
magnitude is \(b=X-D_c\). A low singleton cannot immediately precede or
follow such a trio. If it did, its slot value \(D_r\) would satisfy
\(|D_c-D_r|<6\) (one fixed trio step 2 and the neighboring step strictly
below 4). Then \[
\frac{a-b}{2}=X+\frac{D_c-D_r}{2}
\] is a positive first-six signed-support form inside \((X-3,X+3)\),
impossible. It is an allowed support form because \(a,b\) are oriented
full first-six words and halving their difference gives coefficients in
\(\{-1,0,1\}\). Its support is nonempty by the distinct numerical bands.

Now if \(q=50\) or 60 had both seventh signs minus, its position \(q+1\)
would be a positive all-plus singleton \(D_{q+1}=3X-(w+Y)\), with
\(w+Y>0\). It is low and immediately precedes the first member at
\(q+2\) of the already forced two-plus trio whose center is \(q+3=53\)
or 63. This contradicts the preceding observation. Hence \textbf{all
four prescribed pairs are both plus}, with opposite first-six words.
Their physical gaps \(g_{158},g_{161},g_{168},g_{171}\) are all exactly
1, as each equals \((D_{q+1}-D_q)/2\).

In particular the four centers have corresponding first-six positive
full magnitudes \(b_q=X-D_q\). Since \(D_{q-1}=D_q-2\) and
\(D_{q-1}\ge2q-3\), one has \(D_q\ge2q-1\) (the same general lower
envelope). Therefore \[
b_{50}\in[48,50),\quad b_{53}\in[42,44),\quad
b_{60}\in[28,30),\quad b_{63}\in[22,24).
\]

\subsection{3. One high singleton and one two-plus trio between the
forced
blocks}\label{one-high-singleton-and-one-two-plus-trio-between-the-forced-blocks}

Slots 55--58 are strictly between the complete trios 52--54 and 59--61,
so they consist either of four singletons or one complete trio and a
singleton. Any singleton at 55 or 58 immediately abuts a two-plus trio
and therefore must be high. Consecutive singletons of the same side are
impossible: their positive first-six magnitudes \(a,a'\) would differ by
the slot step, which is smaller than 4, whereas distinct full first-six
magnitudes differ by at least 6. Thus four consecutive singleton slots
would alternate high/low, and their endpoints could not both be high.
Exactly one singleton and one trio occur, and the singleton must be at
an endpoint, either 55 or 58. In particular it is high. With
\(S_6=\sum_{i=1}^6x_i\), its first-six magnitude is \(a=3X+D_r\),
implying the necessary affine root row \[
\boxed{S_6\ge3X+109}
\] since \(r\ge55\) and \(D_{55}\ge109\).

The remaining trio must be two-plus. An adjacent one-plus trio with
center \(D_c\) would correspond to the positive first-six magnitude
\(b=X+D_c\), and the high singleton would give \((a-b)/2=X+(D_r-D_c)/2\)
in the same impossible middle forbidden window. Therefore there is a
fifth distinct \emph{small} full first-six magnitude \(X-D_c\): it is in
\([36,38)\) if the trio center is 56, or \([34,36)\) if the center is
57.

One further exact signed-support relation useful later: \[
t_1=(b_{50}-b_{53})/2,\qquad t_2=(b_{60}-b_{63})/2
\] both lie in \([3,4)\) (two exact trio steps and one inter-trio step
in \([2,4)\)). If their signed support words are distinct, they must
have opposite signs on some shared coordinate: otherwise their
difference is a nonzero \(\{-1,0,1\}^6\) form with absolute value below
1, contradicting first-six separation at least 3. No extra claim is made
from the separate high-band avoidance: since the first-six forbidden
pair-sum window is only \((2X-2,2X+2)\), its matching half-forms can
give merely a slot difference at least 4, already following from two
rank steps in the applications above.

\subsection{4. Consequences for central
orientations}\label{consequences-for-central-orientations}

The central positive top signs must belong to \(+++\), \(++-\), \(+--\),
\(---\): a trio containing \(D_1\) begins there, so contributes only its
first pattern. Let \(P,N\) be the sums of the positive and negative
sides of the central first-six signs, respectively, and let
\(g=D_1\in[1,2)\). Expanding the displayed total lower bound by the
signed central equation yields:

\begin{longtable}[]{@{}ll@{}}
\toprule\noalign{}
Central top signs & Necessary condition \\
\midrule\noalign{}
\endhead
\bottomrule\noalign{}
\endlastfoot
\(++-\) & \(P\ge X+(111+g)/2\ge X+56\) \\
\(+--\) & \(N\ge X+(107-g)/2>X+52.5\) \\
\(+++\) & \(P\ge(109+g)/2\ge55\) \\
\(---\) & \(N\ge(109-g)/2>53.5\) \\
\end{longtable}

Thus for the fixed earlier pilot orientation with word
\(\texttt{-----+++-}\), the only positive sign in the first six is
\(x_6\), contradicting the \(++-\) row. Equivalently its central
identity gives \(S_6=2x_6+X-g-2<3X\). This independently clears that
central orientation's \textbf{close branch}, and in particular all
earlier live pairs of it that force the close branch. It does not clear
the wide branch, or any whole orientation.

The additional LP screen tests only \(S_6\ge3X+109\) over the previous
close root relaxations. The explicitly listed structural identities are
written mathematics, not an LP enumeration or Lean formalization.

\section{Exact close-triple total-row
screen}\label{exact-close-triple-total-row-screen}

The nine deterministic guarded batches completed all 69 previously
unresolved close-branch duties. The final Fraction-only verifier replays
their full prior 244-duty partition and certifies this result:

\begin{longtable}[]{@{}lrr@{}}
\toprule\noalign{}
Case scope & New exact exclusions & Not excluded by this screen \\
\midrule\noalign{}
\endhead
\bottomrule\noalign{}
\endlastfoot
Close triple \((x_7,x_8,x_9)=(X-1,X,X+1)\) & 20 & 49 \\
Wide terminal pair & Not tested & 122 inherited \\
\end{longtable}

Accounting for the earlier 53 close-only exact certificates, the new
close certificates leave 49 unresolved close and 122 unresolved wide
duties, a total of 171 duties across 122 central orientations.
\textbf{No new whole central orientation is excluded.} No lower bound of
150 or complete floor-148 exclusion is claimed. The 49 were only
numerical survivors of this relaxation, not feasible configurations or
certified full sign chambers.

All 20 rational certificates use the strong written row
\(\sum_{i=1}^6 x_i\ge3X+109\) with a positive multiplier, and all have
weighted RHS \(-1\). The top-three sign pattern among these exact closes
is +-- (10), ++- (8), +++ (1), and --- (1). The earlier analytically
excluded close word -----+++- is included in the 20 (not added to the
count); its exact vector has new-row multiplier \(1/118\) and does not
use a strict inherited row.

Nine serial guard receipts are included. Child elapsed times sum to 6.08
s; maximum child elapsed time was 0.93 s; maximum audited child peak RSS
was 60,840 KiB. All guards passed. No wide or q-slot follow-up was run.
The written theorem on which the new row depends is in
close-triple-slot-grammar.md and is not formalized in Lean.

Standard-library exact replay:

\begin{verbatim}
sha256sum -c SHA256SUMS
python3 check-complete-close-total.py
\end{verbatim}

\section{Provenance and scope}\label{provenance-and-scope}

Correspondence:
\href{mailto:aafimalek2032@gmail.com}{\nolinkurl{aafimalek2032@gmail.com}}.
The proof, computation, and writing in this research note were developed
with AI assistance.

This note uses the earlier
\href{https://github.com/Aafimalek/novel-discoveries/tree/c582f2587588eeeb195048a74552fe80bdc1d0e7/nine-separated-subset-sums/verification/rank-cuts-below-150/only-floor-148}{floor-148
reduction},
\href{https://github.com/Aafimalek/novel-discoveries/tree/a49d954ffe00669ed70ac7cdd68ef4091263abbc/nine-separated-subset-sums/verification/rank-cuts-below-150/short-gap-normal-form}{short-gap
normalization}, and
\href{https://github.com/Aafimalek/novel-discoveries/tree/63d79a72228169e68bca6f18ec54c97a5147422d/nine-separated-subset-sums/verification/rank-cuts-below-150}{central
rank constraints}. The exact certificates are in the accompanying
archive. The twenty exclusions apply to the normalized close branch
only. No improvement to the end-to-end Lean-checked lower bound of 149
is claimed.

\end{document}
